\documentclass[9pt,twocolumn,a4paper]{extarticle}

% ELECO 2026 yazım düzeni: A4, iki sütun, 9 punto ana metin.
\usepackage[english,turkish]{babel}
\usepackage{amsmath}
\usepackage{newtxmath}
\usepackage{fontspec}
% Tectonic's Snap package may not see user-installed fonts. Local copies
% also allow the same files to be used in an uploaded XeLaTeX project.
\IfFileExists{fonts/Times.TTF}{%
  \defaultfontfeatures{Path=fonts/,Extension=.TTF,
    UprightFont=Times,BoldFont=Timesbd,
    ItalicFont=Timesi,BoldItalicFont=Timesbi}%
}{}
\setmainfont{Times New Roman}
\setsansfont{Times New Roman}
\setmonofont{Times New Roman}
\setmathrm{Times New Roman}
\usepackage{indentfirst}
\usepackage{geometry}
\usepackage{titlesec}
\usepackage{caption}
\usepackage{booktabs}
\usepackage{siunitx}
\usepackage{graphicx}
\usepackage{pgfplots}
\usepackage{tikz}
\usepackage{upgreek}
\usepackage{url}
\usepackage[hidelinks]{hyperref}
\usepackage{enumitem}
\usepackage{multicol}
\setlength{\multicolsep}{0pt}
\DeclareMathSizes{9}{9}{7}{7}
\newcommand{\affilmark}{\textsuperscript{\fontsize{7}{7}\selectfont 1}}

\pgfplotsset{compat=1.18,/pgf/number format/use comma}
\sisetup{detect-all,output-decimal-marker={,}}
\urlstyle{same}

% ELECO: ikinci ve sonraki sayfalar üst 2.5 cm; ilk sayfada başlık
% bloğuna ek 5 mm verilerek üst kenar 3 cm yapılır.
\geometry{
  a4paper,
  left=20mm,
  right=20mm,
  top=25mm,
  bottom=37mm,
  columnsep=5mm
}

\pagestyle{empty}
\setlength{\parindent}{0.4mm}
\setlength{\parskip}{0pt}
\renewcommand{\baselinestretch}{1.0}
\setlist[itemize]{nosep,leftmargin=*,topsep=0pt,partopsep=0pt}

% Başlıklar: 10 pt kalın; bölüm ortalı, alt bölüm sola hizalı.
\titleformat{\section}
  {\centering\bfseries\fontsize{10}{10.8}\selectfont}
  {\thesection.}{0.5em}{}
\titleformat{\subsection}
  {\bfseries\fontsize{10}{10.8}\selectfont}
  {\thesubsection.}{0.5em}{}
\titlespacing{\section}{0pt}{9pt}{9pt}
\titlespacing{\subsection}{0pt}{9pt}{9pt}

% Şekil/çizelge başlıkları: 9 pt, iki yana yaslı.
\captionsetup{
  font=normalfont,
  labelfont=bf,
  justification=justified,
  singlelinecheck=false
}
\captionsetup[figure]{name=Şekil,labelsep=colon,position=bottom,skip=6pt}
\captionsetup[table]{name=Çizelge,labelsep=period,position=top,skip=6pt}
\setlength{\textfloatsep}{9pt}
\setlength{\floatsep}{9pt}
\setlength{\intextsep}{9pt}
\AtBeginDocument{%
  \setlength{\abovedisplayskip}{9pt}%
  \setlength{\belowdisplayskip}{9pt}%
  \setlength{\abovedisplayshortskip}{9pt}%
  \setlength{\belowdisplayshortskip}{9pt}%
}

% thebibliography'nin kendi başlığını kaldır; başlık ELECO biçiminde
% numaralı \section{Kaynaklar} olarak ayrıca verilir.
\makeatletter
\renewenvironment{thebibliography}[1]
  {\list{\@biblabel{\@arabic\c@enumiv}}%
     {\settowidth\labelwidth{\@biblabel{#1}}%
      \leftmargin\labelwidth
      \advance\leftmargin\labelsep
      \usecounter{enumiv}%
      \let\p@enumiv\@empty
      \renewcommand\theenumiv{\@arabic\c@enumiv}}%
   \sloppy
   \clubpenalty4000
   \widowpenalty4000
   \setlength{\itemsep}{0pt}%
   \setlength{\parsep}{0pt}%
   \setlength{\parskip}{0pt}}
  {\endlist}
\makeatother

\begin{document}
\selectlanguage{turkish}
\shorthandoff{=}

\twocolumn[
\begin{center}
  \vspace*{5mm}
  {\fontsize{14}{14}\selectfont\bfseries
  Erişimi Kısıtlı Hassas Zamanlama Donanımı için Yüksek Doğruluklu Davranışsal Simülatör:\\
  Rubidyum Osilatörü Üzerinde Bir Vaka Çalışması\par}
  \vspace{2mm}
  {\fontsize{14}{14}\selectfont\bfseries
  A High-Fidelity Behavioral Simulator for Limited-Access Precision Timing Hardware:\\
  A Case Study on a Rubidium Oscillator\par}

  \vspace{9pt}
  {\fontsize{11}{11}\selectfont
  Ömer Faruk Karagöz\affilmark, Ömer Yıldız\affilmark, Batuhan Canlı\affilmark\par}

  \vspace{9pt}
  {\fontsize{10}{10.8}\selectfont
  \affilmark TÜBİTAK BİLGEM MA3, Türkiye\\
  omer.karagoz@std.bogazici.edu.tr, omer.yildiz@tubitak.gov.tr, ahmet.canli@tubitak.gov.tr\par}
  \vspace{18pt}
\end{center}
]

\begin{center}
  {\fontsize{10}{10.8}\selectfont\bfseries Özet\par}
\end{center}
\vspace{-6pt}
{\fontsize{9}{10.8}\selectfont\itshape
\noindent Hassas zamanlama donanımına bağımlı istemci yazılımlarının geliştirilmesi, cihazın maliyeti, sınırlı erişilebilirliği ve uzun durum geçiş süreleri nedeniyle güçleşebilir. Bu çalışmada, ticari düşük gürültülü rubidyum Küresel Navigasyon Sistemi Disiplinli Osilatörü (Global Navigation System Disciplined Oscillator, GNSDO) için yüksek doğruluklu davranışsal bir simülatör sunulmaktadır. Simülatör, cihazın seri arayüzünü, Programlanabilir Cihazlar için Standart Komutlar (Standard Commands for Programmable Instruments, SCPI) protokolünü, durum geçişlerini ve istemciden gözlenebilen zamana bağlı davranışlarını yeniden üretmektedir. Model, üretici kılavuzuna, sınırlı cihaz kayıtlarına ve doğrudan gözlemlere dayanmaktadır. Holdover sırasında zaman aralığı hatası, ısınma ve kilitlenme geçişleri, sağlık bayrakları ve ölçüm değişimleri modellenmiştir. Doğrulama, ağırlıkla kilitli duruma ait sınırlı kayıtlarla arayüz ve ilişki karşılaştırmasına dayanır. Holdover ve ısınma eğrileri model çıktılarıdır; zaman hizalı donanım ölçümleriyle doğrulanmamıştır. Otomatik model testleri ve ayrı bir seri izleme istemcisiyle uçtan uca test, yazılımın çalışmasını sınar. Simülatör, gerçek donanıma sürekli erişim gerektirmeden tekrarlanabilir bir test ortamı sağlayabilmektedir.\par}

\vspace{9pt}
\begin{center}
  {\fontsize{10}{10.8}\selectfont\bfseries Abstract\par}
\end{center}
\vspace{-6pt}
{\selectlanguage{english}\fontsize{9}{10.8}\selectfont\itshape
\noindent Developing client software for precision timing hardware is difficult because of device costs, limited availability, and long state transition times. This paper presents a high-fidelity behavioral simulator for a commercial low-noise rubidium Global Navigation System Disciplined Oscillator (GNSDO). The simulator reproduces the device's serial interface, Standard Commands for Programmable Instruments (SCPI) protocol, state transitions, and time-dependent behavior observable by client software. The model uses the vendor manual, limited device recordings, and direct observations. Time-interval error during holdover, warm-up and lock transitions, health flags, and measurement variations are modeled. Hardware comparison is limited to recorded responses and relationships, mainly in the locked state. Holdover and warm-up curves are model outputs, not time-aligned hardware validation. Automated model tests and an end-to-end test with a separate serial client verify software behavior. The simulator can provide a repeatable test environment without continuous access to the physical device.\par}
\selectlanguage{turkish}

\section{Giriş}
Küresel Navigasyon Uydu Sistemi (\emph{Global Navigation Satellite System}, GNSS) ile disipline edilen osilatörler, telekomünikasyon, ölçüm, enerji ve savunma uygulamalarında izlenebilir zaman ve frekans üretmek için kullanılmaktadır. Bu cihazlar, GNSS sisteminin uzun dönem doğruluğunu yerel bir osilatörün kısa dönem kararlılığıyla birleştirir. İncelenen cihazda yerel referans rubidyum tabanlı bir atomik osilatördür ve cihazın izleme, yapılandırma ve durum bilgilerine seri arayüz üzerinden, Programlanabilir Cihazlar için Standart Komutlar (\emph{Standard Commands for Programmable Instruments}, SCPI) kullanılarak erişilebilmektedir \cite{jlt2015,scpi1999}.

Bu tür bir cihaza bağımlı istemci yazılımlarının geliştirilmesinde fiziksel erişim önemli bir darboğaz oluşturabilir. Donanımın pahalı veya sınırlı sayıda olması, aynı cihazın birden fazla geliştirici ya da sürekli entegrasyon ortamı tarafından eşzamanlı kullanılamaması ve GNSS kaybı, referans kaybı sonrasındaki serbest çalışma (\emph{holdover}) veya ilk açılış gibi durumların kontrollü olarak tekrar edilmesinin zaman alması test kapsamını sınırlar. Gerçek donanımın test ortamında bulunmaması veya kullanımının zor olması, yazılım testinde alternatif uygulamalar ve donanım simülatörleri kullanılmasının temel gerekçelerinden biridir \cite{meszaros2007,morita2017}.

Bu çalışmada, üreticinin \emph{Global Navigation System Disciplined Oscillator} (GNSDO) olarak adlandırdığı Jackson Labs Low Noise Rubidium cihazı için, istemci yazılımının gördüğü seri protokolü ve zamana bağlı davranışı yeniden üreten bir \emph{high-fidelity behavioral simulator} geliştirilmiştir. Buradaki ``high-fidelity'' ifadesi bir tasarım hedefidir; deneysel destek, incelenen komut/yanıt biçimleri ve kayıtlardaki değişken ilişkileriyle sınırlıdır. Zaman alanı modelleri istemci senaryoları üretir; fiziksel cihazın geçici rejimlerini veya metrolojik doğruluğunu nicel olarak yeniden ürettiği gösterilmemiştir. Simülasyon doğrulamasında modelin geçerliliğinin kullanım amacıyla birlikte ele alınması gerektiği literatürde de vurgulanmaktadır \cite{sargent2013}.

Çalışmanın temel katkıları şunlardır:
\begin{itemize}
    \item Gerçek seri aygıtla aynı erişim biçimini sağlayan sanal seri port, SCPI ayrıştırıcı, tekil cihaz-durumu modeli ve zamana bağlı davranış katmanlarından oluşan modüler bir mimari;
    \item holdover, ısınma, kilitlenme, sağlık bayrakları ve ölçüm değişkenleri için tekrarlanabilir davranış modelleri;
    \item donanıma sınırlı erişim koşullarında uygulanabilen, kaynak izlenebilirliği, kayıt karşılaştırması, ilişki kontrolleri ve fiziksel değişmez testlerini birleştiren doğrulama yaklaşımı.
\end{itemize}

\section{Arka Plan ve İlgili Çalışmalar}
\subsection{Referans Cihaz ve Komut Arayüzü}
Referans cihaz Jackson Labs Technologies Low Noise Rubidium GNSDO'dur \cite{jlt2015}. Cihazın kimlik yanıtındaki GPSDO, Küresel Konumlama Sistemi ile Disipline Edilen Osilatör (\emph{Global Positioning System Disciplined Oscillator}) anlamındadır. Komut adlarındaki CSAC, çip ölçekli atomik saat (\emph{Chip-Scale Atomic Clock}) kısaltmasıdır; bu arayüz adı, incelenen rubidyum varyantının tüm CSAC donanım özelliklerini taşıdığı anlamına gelmez. Üretici, ürün adında GNSDO kısaltmasını kullanmaktadır; bu nedenle makalede ürün adı için GNSDO, genel teknoloji sınıfı için ise GNSS ile disipline edilen osilatör ifadesi kullanılmaktadır. Üretici kılavuzuna göre cihazın RS-232 seri arayüzünün fabrika ayarı 115200 baud hız, 8 veri biti, parite yok, 1 stop biti ve akış kontrolü yok biçimindedir; bu yapı kısaca 8N1 olarak gösterilmektedir. Komut arayüzü SCPI sözdizimini kullanır; SCPI standardı programlanabilir ölçüm cihazları için donanım arayüzünden bağımsız bir komut dili tanımlar \cite{scpi1999}.

Cihazda rubidyum osilatörü GNSS veya harici saniyede bir darbe (\emph{one pulse per second}, 1PPS) referansına kilitleyen ve filtre osilatörünü rubidyum referansına kilitleyen iki ayrı servo döngüsü bulunmaktadır \cite{jlt2015}. \texttt{SYNC:TINT?} komutu, seçili döngü için \emph{Time Interval} (TINT) olarak adlandırılan 1PPS faz farkını raporlar. \texttt{SYNC:FEEstimate?} komutu ise üreticinin 1000~s ölçüm aralığında Allan varyansına benzer bir büyüklük olarak tanımladığı \emph{Frequency Error Estimate} (FEE; frekans hata kestirimi) değerini döndürür. \texttt{SYNC:HEALTH?} komutu çalışma süresi, holdover ve faz farkı gibi koşullara bağlı sağlık bayraklarını onaltılık bir bit alanı olarak verir. Tanısal komutlarda kullanılan \emph{Electronic Frequency Control} (EFC; elektronik frekans kontrolü) değerleri de iç referans osilatörüne uygulanan kontrol büyüklüğünü ifade eder \cite{jlt2015}.

\subsection{Davranışsal Simülasyon ve Test Amaçlı Donanım İkamesi}
Meszaros'un \emph{Fake Object} sınıfı, gerçek bağımlılığın yerine geçen basitleştirilmiş çalışan uygulamaları kapsar \cite{meszaros2007}. Üretim kayıtlarından test ikamesi oluşturma \cite{tiwari2024} ve donanım arızalarını simüle ederek istemci testi yapma \cite{morita2017}, bu yaklaşımın ilgili örnekleridir. Burada amaç, üçüncü taraf cihazın firmware'ine erişmeden gözlenebilir seri davranışı yeniden kurmaktır.

\subsection{Simülasyon Doğrulaması}
Sargent, kavramsal model, uygulama, veri ve operasyonel geçerliliği ayırır; doğrulama kullanım amacına bağlıdır \cite{sargent2013}. Burada kayıt karşılaştırmaları, parametre izlenebilirliği ve yazılım testleri ayrı kanıt türleri olarak raporlanmıştır.

\subsection{Saat ve Holdover Modellemesi}
Saat fazı ve frekans hatası arasındaki ilişki klasik saat modellerinin temelidir. Zucca ve Tavella, saat hatasını stokastik diferansiyel denklemlerle ifade ederek frekans süreçleri ile Allan varyansı arasındaki ilişkiyi ayrıntılı biçimde incelemiştir \cite{zucca2005}. Bu çalışmanın amacı osilatörün tam stokastik spektrumunu modellemek değildir. Bunun yerine, istemci testleri için gerekli zaman alanı davranışları daha düşük karmaşıklıkta ve deterministik biçimde temsil edilmektedir.

\section{Tasarım Hedefleri ve Kapsam}
Simülatörün temel hedefi, gerçek cihazla çalışan istemci yazılımında değişiklik gerektirmeden aynı bağlantı ve komut modelini sunmaktır. Bu amaçla aşağıdaki doğruluk boyutları esas alınmıştır:
\begin{enumerate}
    \item \textbf{Arayüz doğruluğu:} Komut adları, kısa/uzun SCPI biçimleri, satır sonlandırma ve çok satırlı yanıtların biçimi;
    \item \textbf{Durum tutarlılığı:} Özet sorgular ile tekil sorguların aynı cihaz durumundan türetilmesi;
    \item \textbf{Zamansal davranış hedefi:} Kılavuz ve açık varsayımlara dayalı ısınma, kilitlenme, holdover ve sağlık geçişlerinin üretilebilmesi;
    \item \textbf{İlişkisel doğruluk:} Birbirine bağlı ölçümlerin kayıtlarda ve kılavuzda gözlenen ilişkileri koruması;
    \item \textbf{Tekrarlanabilirlik:} Aynı başlangıç koşulları altında aynı test çalışmasının aynı çıktıları üretmesi.
\end{enumerate}

Spektral faz gürültüsü, Allan sapması ve cihazlar arası varyasyon kapsam dışıdır. Kayıtlı \texttt{HELP?} listesindeki 66 cihaz sorgusunun 50'si desteklenir; \texttt{HELP?} sayım dışında tutulmuştur. Bu sayı işlev kapsamıdır; tüm kısa/uzun yazım birleşimlerinin veya ayar komutlarının doğruluğunu göstermez. Desteklenmeyen 16 sorgu ve etkileri Çizelge~\ref{tab:unsupported}'te verilmiştir.

\begin{table}[t]
\caption{Desteklenmeyen sorgular ve istemciye etkileri}
\label{tab:unsupported}
\centering
\fontsize{9}{10.8}\selectfont
\setlength{\tabcolsep}{3pt}
\begin{tabular}{p{0.46\columnwidth}p{0.46\columnwidth}}
\toprule
\textbf{Sorgu grubu} & \textbf{Kapsam sınırı} \\
\midrule
\texttt{MAC:} altında \texttt{RS232?}, \texttt{STEER?}, \texttt{STATUS?}, \texttt{MODE?}, \texttt{DDSCENTER?}, \texttt{TECCONTROL?}, \texttt{TCXO?}, \texttt{SIGNAL?}, \texttt{HEATPACKAGE?}, \texttt{TEMPERATURE?}, \texttt{SN?}, \texttt{FWREV?}, \texttt{LIFETIME?} (13)
 & MAC alt sorgularını kullanan modül izleme/tanı istemcileri çalışmaz. DDS merkezi dışında karşılık gelen alanlar \texttt{CSAC:*} sorgularıyla; DDS merkezi ise toplu \texttt{CSAC?}/\texttt{MAC?} yanıtıyla okunabilir. \\
\texttt{DIAG:ROSC:EFC:} altında \texttt{REL:CSAC?}, \texttt{REL:FILTER?}, \texttt{ABS:FILTER?} (3)
 & Döngüye özel bağıl kontrol ve filtre mutlak kontrol izlemesi eksiktir; genel EFC sorguları bu işlevlerin tamamının yerini tutmaz. \\
\bottomrule
\end{tabular}
\end{table}

\section{Simülatör Mimarisi}
Simülatör dört katmandan oluşmaktadır:
\begin{itemize}
    \item \textbf{Sanal seri arayüz:} \texttt{socat} ile oluşturulan sözde terminal (\emph{pseudo-terminal}) üzerinden istemciye standart bir \texttt{/dev/tty*} aygıtı sunulur.
    \item \textbf{SCPI ayrıştırıcı:} Komutlar kısa ve uzun biçimleriyle normalize edilir, uygun işleyiciye yönlendirilir ve desteklenmeyen girişlerde cihazla uyumlu hata yanıtı üretilir.
    \item \textbf{Cihaz-durumu katmanı:} Kimlik, yapılandırma, servo durumu, ölçümler ve sağlık bilgileri tek bir durum nesnesinde tutulur. Böylece birbiriyle ilişkili sorgular aynı anda çelişkili değer üretmez.
    \item \textbf{Davranış modelleri:} Zamana bağlı fiziksel ve mantıksal davranışlar giriş/çıkış katmanından ayrılmış saf fonksiyonlar olarak uygulanır.
\end{itemize}

Davranış modellerinin geçen süreyi doğrudan parametre olarak alması, saatler süren gerçek zamanı beklemeden uzun holdover veya ilk açılış senaryolarının birim testinde hızlandırılmış biçimde sınanabilmesini sağlar.

\section{Davranış Modelleri}
\subsection{Holdover Sırasında Zaman Aralığı Hatası}
GNSS referansı kaybedildiğinde servo döngüsü referansa göre düzeltme yapamaz ve yerel osilatörün artık frekans hatası faz hatasına dönüşür. Kısa ve orta süreli istemci testleri için birinci dereceden yaklaşım kullanılmıştır:
\begin{equation}
    x(t)=x_0+y_h t,
    \label{eq:holdover}
\end{equation}
burada $t$ holdover başlangıcından geçen süre (s), $x(t)$ ve $x_0$ sırasıyla zaman aralığı hatası ve başlangıç faz farkı (s), $y_h$ ise işaretli, boyutsuz etkin fraksiyonel frekans hatasıdır.

Üretici kılavuzu \texttt{SYNC:FEEstimate?} değerini 1000~s aralığında Allan varyansına benzer bir büyüklük olarak tanımladığından, bu değer doğrudan işaretli fiziksel frekans sapması olarak yorumlanmamıştır. Simülatörde $|y_h|$ için kaydedilmiş FEE büyüklükleri başlangıç referansı olarak kullanılmakta; işaret ve ölçekleme açık bir model parametresi olarak tutulmaktadır. Bu yaklaşım, kaynak tarafından desteklenmeyen bir fiziksel yorumu modele gömmek yerine varsayımı görünür kılar.

Bu çalışmanın sabitlenmiş kod sürümünde ikinci dereceden yaşlanma terimi kullanılmaz. Kılavuzdaki $8\times10^{-14}$ Allan sapması değerinden yaşlanma katsayısı çıkarılmamıştır. \texttt{holdover\_frequency\_error} parametresi $y_h$ değerini açıkça belirler; ayarlanmazsa girişteki FEE pozitif eğim olarak eşlenir ve bu eşleme bir model varsayımıdır. Kod sürümü ve yeniden üretim yöntemi Bölüm~\ref{sec:availability}'te verilmiştir.

\subsection{Deterministik Ölçüm Değişimi}
Gerçek cihazda sıcaklık, TINT ve benzeri büyüklükler tamamen sabit değildir. Ancak her sorguda bağımsız rastgele sayı üretmek hem sorgu sıklığına bağlı yapay değişim oluşturur hem de regresyon testlerinin tekrar üretilebilirliğini bozar. Bu nedenle gözlenebilir küçük değişimler zamanın deterministik bir fonksiyonu olarak modellenmiştir:
\begin{equation}
 v(t_u)=v_0+A\frac{\sum_{k=0}^{2}w_k\sin\!\left(2\uppi r_k t_u/P+\upvarphi_k\right)}{\sum_{k=0}^{2}w_k},
 \label{eq:variation}
\end{equation}
Burada $t_u$ simülatör başlangıcından geçen süre (s), $v(t_u)$ çıktı, $v_0$ nominal değer ve $A$ sapma sınırıdır; son üç büyüklük aynı kanal birimindedir (TINT için s, sıcaklık için $^{\circ}$C). $P$ temel periyot (s), $k=0,1,2$ bileşen indisi, $w_k=(1,1/2,1/4)$ boyutsuz ağırlıklar ve $r_k=(1,\upphi,\upphi^2)$ boyutsuz frekans oranlarıdır. $\upphi=(1+\sqrt{5})/2$ altın oran, $\upvarphi_k=r_k\uptheta_c$ radyan cinsinden sabit fazdır; $\uptheta_c$ kanal adından deterministik üretilir. Normalizasyon sayesinde değişim $\pm A$ aralığıyla sınırlı kalır. Bu yapı fiziksel gürültünün spektral modeli olarak değil, istemci testleri için zaman-tutarlı ve tekrarlanabilir küçük değişim üretmek amacıyla kullanılmaktadır.

\subsection{Isınma ve Kilitlenme}
Baskılı devre kartı (\emph{printed circuit board}, PCB) sıcaklığı için birinci dereceden termal yaklaşım kullanılmıştır:
\begin{equation}
 T(t_w)=T_{\infty}-\left(T_{\infty}-T_0\right)e^{-t_w/\uptau},
 \label{eq:thermal}
\end{equation}
burada $t_w$ soğuk açılıştan geçen süre (s), $T_0$ başlangıç sıcaklığı, $T_{\infty}$ kararlı sıcaklık ($^{\circ}$C) ve $\uptau$ termal zaman sabitidir (s). Varsayılan parametreler, eldeki gözlem ve kayıtlarla uyumlu olacak şekilde $T_0=\SI{25}{\celsius}$, $T_{\infty}\approx\SI{52.8}{\celsius}$ ve $\uptau\approx\SI{90}{s}$ olarak seçilmiştir. $T_0$ ortam koşuluna bağlı olduğu için yapılandırılabilir tutulmuştur; $\uptau$ üretici spesifikasyonu değil, gözlemsel bir model parametresidir.

Üretici kılavuzu, GNSS anteni bağlıyken tam kilit göstergesinin güç verildikten sonra tipik olarak 20 dakikadan kısa sürede oluştuğunu belirtmektedir \cite{jlt2015}. Servo durumları bu zaman ölçeğini koruyacak şekilde ısınma, kilitlenme ve kilitli durumlara ayrılmıştır. Kılavuzdaki ``2 dakikadan kısa sürede atomik kilit'' ifadesi, CSAC kullanılan belirli varyanta ait olduğundan referans rubidyum cihazının doğrulanmış ısınma süresi olarak alınmamıştır. Kodun 120~s atomik kilit geçişi de bu cihaz için ölçülmüş sonuç değil, senaryo varsayımıdır.

\subsection{Sağlık Durumu}
\texttt{SYNC:HEALTH?} için sağlık bayrakları üretici kılavuzundaki koşullara göre oluşturulmaktadır. Modelde özellikle aşağıdaki kurallar uygulanmaktadır: çalışma süresi 200~s'den küçükse \texttt{0x8}; holdover süresi 60~s'yi aşarsa \texttt{0x10}; referansa göre faz farkı 210~ns'yi aşarsa \texttt{0x4}. Bayraklar bit düzeyinde mantıksal VEYA işlemiyle birleştirilir. Örneğin uzun holdover sırasında faz farkının 210~ns'yi aşması \texttt{0x14} sonucunu üretir. Sağlık bayrağındaki 210~ns koşulu ile \texttt{SYNC:TINT:THR} komutunun 220~ns varsayılan ani faz hizalama (\emph{jam-sync}) eşiği ayrı kavramlardır ve modelde birbirine karıştırılmamıştır \cite{jlt2015}.

\section{Modelin İnşası ve Parametre İzlenebilirliği}
Model parametreleri dört kaynak sınıfından biriyle ilişkilendirilmiştir: üretici kılavuzu, kaydedilmiş cihaz çıktısı, doğrudan operatör gözlemi ve açık varsayım. Böylece her sayısal değerin kaynağı ve güven düzeyi incelenebilir hale gelir. Çizelge~\ref{tab:trace} bu ayrımı özetlemektedir.

\begin{table}[t]
\caption{Model parametrelerinin izlenebilirliği}
\label{tab:trace}
\centering
\fontsize{9}{10.8}\selectfont
\setlength{\tabcolsep}{3pt}
\begin{tabular}{p{0.25\columnwidth}p{0.28\columnwidth}p{0.38\columnwidth}}
\toprule
\textbf{Parametre / ilişki} & \textbf{Değer} & \textbf{Kaynak ve durum} \\
\midrule
Seri arayüz & 115200, 8N1 & Kılavuz, doğrudan \\
Sağlık faz koşulu & $|TINT|>210$ ns & Kılavuz, doğrudan \\
Ani faz hizalama (jam-sync) eşiği & 220 ns varsayılan & Kılavuz, doğrudan \\
Holdover bayrağı & $t>60$ s & Kılavuz, doğrudan \\
İlk çalışma bayrağı & $t_u<200$ s & Kılavuz, doğrudan \\
$|y_h|$ başlangıç aralığı & $1{,}3$--$2{,}5\times10^{-11}$ & Kayıt + model eşleme varsayımı \\
Kilitli TINT model sınırı & $\pm12$ ns & 5 kayıtla ölçek seçimi \\
CSAC--PCB sıcaklık farkı & $+1{,}32\,^{\circ}$C & Kayıtlar (4 çift) \\
EFC mutlak/bağıl ilişki & $a=200r$ & Kayıtlar \\
$\uptau$ & yaklaşık 90 s & Operatör gözlemi / model parametresi \\
$T_0$ & $25\,^{\circ}$C & Varsayım, yapılandırılabilir \\
\bottomrule
\end{tabular}
\end{table}

Bu tablo aynı zamanda modelin kanıt düzeyini sınırlayan noktaları açıkça göstermektedir. Özellikle FEE değerinden etkin holdover eğimine yapılan eşleme ve termal zaman sabiti, doğrudan üretici spesifikasyonu olarak sunulmamaktadır.

\section{Doğrulama Yöntemi}
Donanım karşılaştırması ile yazılımın model denklemlerine uygunluğu ayrı değerlendirilmiştir. Kayıtlarda bazı saat/tarih yanıtları vardır; ancak sorguların tamamını zaman hizalı bir holdover veya soğuk açılış serisine bağlayacak zaman damgaları bulunmamaktadır.

\subsection{Komut ve Yanıt Karşılaştırması}
Kaydedilmiş gerçek cihaz çıktıları ile simülatör yanıtları aynı komutlar için karşılaştırılmıştır. Kimlik, seri numarası ve sabit yapılandırma alanlarında tam metin eşleşmesi; çok satırlı komutlarda alan sırası, etiketler ve satır yapısı kontrol edilmiştir.

\subsection{İlişki Tutarlılığı}
Tek tek değerlerin yanı sıra değişkenler arası ilişkiler doğrulanmıştır. Dört EFC çiftinde $a=200r$ bulunmuştur: $r$ yüzde olarak raporlanan sayısal değer, $a$ ise ppt ($10^{-12}$) cinsinden sayısal değerdir. Örneğin $r=-0{,}410000$ (\%) için $a=-82$ (ppt); katsayı $200$'ün birimi ppt/\%'dir. Dört sıcaklık çiftinde CSAC--PCB farkı 1{,}2612--1{,}3881~$^{\circ}$C aralığındadır; modeldeki 1{,}32~$^{\circ}$C ofsetin kayıtlara göre en büyük mutlak farkı 0{,}0681~$^{\circ}$C'dir. Bu, zaman hizalı ölçüm hatası veya cihaz doğruluğu değildir. Simülatör bu ilişkileri ayrı sabit yanıtlar yerine ortak durum üzerinden üretmektedir.

\subsection{Fiziksel ve Mantıksal Değişmezler}
Model için otomatik test edilen başlıca değişmezler; sıcaklığın kararlı değere sınırlı yaklaşması, aynı başlangıç koşullarında aynı serinin üretilmesi, holdover süresinin monoton artması, sağlık bayraklarının koşullarla tutarlı olması ve özet/tekil sorguların aynı cihaz durumunu yansıtmasıdır.

\subsection{Zamansal Davranış Kontrolü}
Model fonksiyonlarına geçen süre verilerek uzun süreli senaryolar beklemeden sınanmıştır. Bu testler denklem ve eşiklerin kodda doğru uygulanmasını kontrol eder. Holdover eğimi, ısınma zaman sabiti, kilitlenme anı ve donanımın sağlık geçiş gecikmeleri zaman hizalı gerçek cihaz verisiyle nicel olarak doğrulanmamıştır; sıcaklığın birkaç dakikada yükselmesi yalnızca nitel operatör gözlemidir.

\section{Sonuçlar}
\subsection{Arayüz ve İlişki Doğruluğu}
Çizelge~\ref{tab:compare}, mevcut kayıt kümesi üzerinde yapılan karşılaştırmaların bir bölümünü göstermektedir.

\begin{table}[t]
\caption{Simülatör çıktısının kaydedilmiş cihaz çıktısıyla karşılaştırılması}
\label{tab:compare}
\centering
\fontsize{9}{10.8}\selectfont
\setlength{\tabcolsep}{3pt}
\begin{tabular}{p{0.26\columnwidth}p{0.44\columnwidth}p{0.20\columnwidth}}
\toprule
\textbf{Sorgu} & \textbf{Kayıt / beklenen ilişki} & \textbf{Sonuç} \\
\midrule
\texttt{*IDN?} & Jackson Labs, LN Rb GPSDO (PRE), Firmware Rev 1.17 & Tam \\
\texttt{SYST:ID:SN?} & 122301668 & Tam \\
\texttt{CSAC:SN?} & 2209MX04906 & Tam \\
\texttt{CSAC:FW?} & V1.0.23 & Tam \\
\texttt{MEAS?} & dört etiketli satır & Biçim tam \\
\texttt{DIAG?} & $a=200r$ (ppt/\%) & 4/4 \\
CSAC--PCB & $+1{,}32\,^{\circ}$C ofset & En çok 0{,}0681~$^{\circ}$C fark \\
\texttt{SYNC:FEE?} & $1{,}3$--$2{,}5\times10^{-11}$ & Aralık \\
\texttt{SYNC:TINT?} & 5 kayıt: $[-11{,}79;\ 11{,}33]$ ns & Gözlenen aralık \\
\bottomrule
\end{tabular}
\end{table}

Çizelgedeki tam eşleşmeler yalnızca listelenen sabit alanlara aittir. Beş TINT örneği zaman serisi veya hata referansı değildir; yaklaşık $\pm12$~ns model değişim sınırını gerekçelendirir, $\pm12$~ns ölçüm doğruluğu iddiası oluşturmaz. Kaynakta kilitli ortalama faz doğruluğu için verilen değer de bu anlık değişim aralığıyla aynı ölçüt değildir.

\subsection{Holdover Davranışı}
Örnek bir $|y_h|=1{,}8\times10^{-11}$ değeri için Denklem~\ref{eq:holdover}, $x_0=0$ koşulunda bir saat sonunda 64{,}8~ns ve bir gün sonunda yaklaşık \SI{1.56}{\micro\second} zaman aralığı hatası üretmektedir. 210~ns sağlık koşulu yaklaşık 3{,}24~saatte aşılmaktadır. Şekil~\ref{fig:holdover}, revize kodun saf holdover fonksiyonundan üretilmiştir; ölçüm değişimi eklenmemiştir. Grafikte saat, hesapta $t=3600h$ ile saniye kullanılır.

\begin{figure}[t]
\centering
\begin{tikzpicture}
\begin{axis}[
 width=\columnwidth,
 height=0.62\columnwidth,
 xlabel={Holdover süresi (saat)},
 ylabel={Zaman aralığı hatası (ns)},
 xmin=0,xmax=12,
 ymin=0,ymax=800,
 grid=both,
 tick label style={font=\scriptsize},
 label style={font=\small},
]
\addplot[thick] table[x=hours,y=tint_ns,col sep=comma] {data/holdover.csv};
\addplot[dashed,domain=0:12,samples=2] {210};
\addplot[only marks,mark=*] coordinates {(3.24074,210)};
\node[font=\scriptsize,anchor=south west] at (axis cs:3.35,210) {3,24 saat};
\node[font=\scriptsize,anchor=south east] at (axis cs:11.8,210) {210 ns sağlık koşulu};
\end{axis}
\end{tikzpicture}
\caption{Model çıktısı: $x_0=0$, $y_h=1{,}8\times10^{-11}$ için holdover eğrisi; donanım ölçümü değildir}
\label{fig:holdover}
\end{figure}

Bu hesap bir metrolojik holdover performans garantisi değildir; yalnızca simülatörde kullanılan etkin frekans hata parametresinin istemci tarafından gözlenen TINT değerine nasıl yansıtıldığını gösterir.

% The final page uses multicol's balanced columns. Its figure is placed
% inline so that its height participates in the column balancing.
\onecolumn
\begin{multicols}{2}
\subsection{Sağlık Durumu Geçişleri}
Holdover 60~s'yi aştığında \texttt{0x10} bayrağı etkinleşir. Faz farkı 210~ns'yi aştığında \texttt{0x4} eklenir ve birleşik durum \texttt{0x14} olur. Güç verildikten sonraki ilk 200~s için \texttt{0x8} bayrağı ayrıca modellenmektedir. Bu kurallar üretici kılavuzundaki sağlık bitleriyle doğrudan uyumludur \cite{jlt2015}.

\subsection{Isınma ve Kilitlenme}
Denklem~\ref{eq:thermal}'deki varsayılan parametrelerle PCB sıcaklığı yaklaşık birkaç dakika içinde kararlı değere yaklaşırken, GNSS kilit durumu daha yavaş ilerlemekte ve modelde yaklaşık 20 dakikalık zaman ölçeğinde etkinleşmektedir. Şekil~\ref{fig:warmup} saf \texttt{thermal\_ramp} fonksiyonundan üretilmiştir; ölçüm değişimi eklenmemiştir. Yatay eksendeki dakika değeri hesapta $t_w=60m$ ile saniyeye çevrilir; kesikli çizgi ölçülmüş kilitlenme anı değildir.

\begin{center}
\begin{minipage}{\columnwidth}
\centering
\begin{tikzpicture}
\begin{axis}[
 width=\columnwidth,
 height=0.62\columnwidth,
 xlabel={Açılıştan itibaren süre (dk)},
 ylabel={Kart sıcaklığı ($^{\circ}$C)},
 xmin=0,xmax=25,
 ymin=24,ymax=55,
 grid=both,
 tick label style={font=\scriptsize},
 label style={font=\small},
]
\addplot[thick] table[x=minutes,y=pcb_c,col sep=comma] {data/warmup.csv};
\addplot[dashed] coordinates {(20,24) (20,55)};
\end{axis}
\end{tikzpicture}
\captionof{figure}{Model çıktısı: PCB ısınması ve 20~dk senaryo kilit eşiği; donanım zaman serisi değildir}
\label{fig:warmup}
\end{minipage}
\end{center}

Saf model fonksiyonları aynı geçen süre ve parametrelerle aynı sonucu verir. Gerçek zamanlı seri sorguların örnekleme anları değişebileceğinden bu özellik, farklı canlı koşuların bayt düzeyinde eşitliği anlamına gelmez.

\subsection{Seri İstemciyle Uçtan Uca Örnek}
Revizyon için hazırlanan küçük Python/pySerial izleme istemcisi, simülatör modüllerini içe aktarmadan ayrı süreçte çalıştırılmıştır. \texttt{socat} iki sanal seri portu bağlamış; istemci 115200 baud, 8N1 ve CR/LF ile kimlik, kilit, holdover, TINT ve sağlık sorgularını ayrıştırmıştır. Normal durumda kilit 1 iken \texttt{SYNC:HOLD:INIT} sonrasında kilit 0 ve holdover alarmı 1 olmuş; gerçek zamanda 62~s bekleme sonunda \texttt{0x10} süre bayrağı görülmüştür. \texttt{SYNC:HOLD:REC:INIT} sonrasında kilit 1'e dönmüş ve holdover alarmı temizlenmiştir. İlk 200~s için ayrıca \texttt{0x8} bulunduğundan sağlık alanının tamamı yerine ilgili bit sınanmıştır. Koşu kaydı ve istemci kaynak kodu sürümle birlikte sunulmuştur. Bu test seri taşıma, komut ayrıştırma ve alarm akışını gösterir; mevcut bir saha uygulamasıyla veya fiziksel cihazla yapılmış kabul testi değildir.

\section{Sınırlamalar}
Yüksek doğruluk hedefinin kanıtı, incelenen arayüz alanları ve kayıtlı ilişkilerle sınırlıdır; tüm 66 sorgu, geçici rejimler veya fiziksel faz doğruluğu için genellenemez. Faz gürültüsü, Allan sapması ve osilatörün ayrıntılı stokastik karakteri modellenmemektedir. Ölçüm değişimleri test tekrar edilebilirliği için bilinçli olarak deterministiktir; bu nedenle simülatör osilatör kararlılığının metrolojik değerlendirilmesinde kullanılamaz.

Doğrulama tek bir fiziksel cihaz örneği ve sınırlı kayıt kümesiyle yapılmıştır. Bazı kayıtlarda zaman damgası bulunmadığından tüm davranışlar zaman hizalı gerçek/simülatör eğrileriyle karşılaştırılamamıştır. Ayrıca FEE değerinden etkin holdover eğimine yapılan eşleme üretici kılavuzunun doğrudan tanımladığı bir fiziksel ilişki değildir ve model parametresi olarak açıkça tutulmaktadır. Termal zaman sabiti de üretici spesifikasyonu değil, gözleme dayalı bir parametredir.

\section{Sonuç ve Gelecek Çalışmalar}
Bu çalışmada seri/SCPI arayüzünün desteklenen bölümünü, ortak cihaz durumunu ve açık varsayımlara dayalı zamansal senaryoları üreten bir simülatör geliştirilmiştir. Simülatör, gerçek donanımın sürekli olarak test ortamında bulunmasına gerek kalmadan istemci yazılımının aynı komut arayüzü üzerinden çalıştırılmasına ve uzun süreli durumların tekrarlanabilir biçimde sınanmasına olanak vermektedir.

Sınırlı kayıtlar arayüz ve ilişki kontrollerini destekler; zamansal donanım doğrulamasının yerini tutmaz. Gelecek çalışma, zaman damgalı holdover/soğuk açılış kayıtları ve mevcut saha istemcileriyle karşılaştırmalı testlerdir.

\section{Erişilebilirlik ve Kod Sürümü}
\label{sec:availability}
Simülatör kodu ve deney kayıtları \url{https://github.com/Pangea567/gnsdo-scpi-simulator} deposunda \nolinkurl{eleco-2026-review-v1} etiketi altında bulunmaktadır. Bildiri dosyaları \path{papers/eleco/} dizinindedir. Bu dizindeki \path{generate_evidence.py}, Şekil~\ref{fig:holdover}--\ref{fig:warmup} verilerini üretir; veriler \path{data/holdover.csv} ve \path{data/warmup.csv} dosyalarındadır. \path{data/manifest.json} kaynak commit kimliğini, dosya SHA-256 özetlerini ve parametreleri içerir. Yeniden üretim komutları \path{REPRODUCIBILITY.md}, 66 sorgunun destek durumu \path{data/query_coverage.csv} dosyasında verilmiştir.

\section{Kaynaklar}
\begin{thebibliography}{99}

\bibitem{jlt2015}
Jackson Labs Technologies, Inc., ``Low Noise Rubidium GNSDO User Manual,'' Document 80200526, Version 1.3, Aug. 2015.

\bibitem{scpi1999}
SCPI Consortium, \emph{1999 SCPI Syntax \& Style}. SCPI Consortium / IVI Foundation, 1999.

\bibitem{meszaros2007}
G. Meszaros, \emph{xUnit Test Patterns: Refactoring Test Code}. Addison-Wesley Professional, 2007.

\bibitem{morita2017}
K. Morita, M. Sato, and S. Wakui, ``Development of the Support System for a Software Testing Using a Hardware Failure Simulation,'' \emph{Journal of the Japan Society for Precision Engineering}, vol. 83, no. 1, pp. 101--107, 2017, doi: 10.2493/jjspe.83.101.

\bibitem{sargent2013}
R. G. Sargent, ``Verification and Validation of Simulation Models,'' \emph{Journal of Simulation}, vol. 7, no. 1, pp. 12--24, 2013, doi: 10.1057/jos.2012.20.

\bibitem{tiwari2024}
D. Tiwari, M. Monperrus, and B. Baudry, ``Mimicking Production Behavior With Generated Mocks,'' \emph{IEEE Transactions on Software Engineering}, vol. 50, no. 11, pp. 2921--2946, 2024, doi: 10.1109/TSE.2024.3458448.

\bibitem{zucca2005}
C. Zucca and P. Tavella, ``The Clock Model and Its Relationship with the Allan and Related Variances,'' \emph{IEEE Transactions on Ultrasonics, Ferroelectrics, and Frequency Control}, vol. 52, no. 2, pp. 289--296, 2005, doi: 10.1109/TUFFC.2005.1406554.

\end{thebibliography}

\end{multicols}
\end{document}
